\section*{अध्याय \ref{ch:b2:geological-hazards} --- भूवैज्ञानिक संकट और जोखिम}

\begin{solution}{exo:b2:geological-hazards:1}
(क) $1 - e^{-30/100} = 0.259$। (ख) $e^{-100/100} = e^{-1} = 0.368$: तीन में से एक से अधिक शताब्दी
उसके बिना बीत जाती है। (ग) $T = -50/\ln 0.95 = 975$ वर्ष।
\end{solution}

\begin{solution}{exo:b2:geological-hazards:2}
(क) कस्बा A: $(1/200) \times 5000 \times 0.02 = 0.5$ भवन प्रति वर्ष; कस्बा B: $(1/50) \times 800
\times 0.10 = 1.6$। (ख) कस्बा B, जिसके उद्भासित भवन 6.25 गुना कम हैं। (ग) उसका संकट 4 गुना
अधिक बारंबार है और उसकी सुभेद्यता 5 गुना अधिक: मिलकर 20 का गुणक, जो
कम \omterm{def:b2:geological-hazards:risk}{उद्भासन} पर भारी पड़ता है।
\end{solution}

\begin{solution}{exo:b2:geological-hazards:3}
(क) $\sqrt{9.81 \times 4000} = \qty{198}{m.s^{-1}}$; $\sqrt{9.81 \times 200} = \qty{44.3}{m.s^{-1}}$;
$\sqrt{9.81 \times 10} = \qty{9.90}{m.s^{-1}}$। (ख) $\num{1.5e6}/198.1 = \qty{7570}{s}$, \qty{2.10}{h}।
(ग) $800/3.6 = \qty{222.2}{m.s^{-1}}$, और $h = c^2/g = 222.2^2/9.81 = \qty{5030}{m}$।
\end{solution}

\begin{solution}{exo:b2:geological-hazards:4}
(क) $\tan\ang{33}/\tan\ang{25} = 0.6494/0.4663 = 1.39$। (ख) $F = 1$ जब $\beta = \varphi' =
\ang{33}$। (ग) $c' = 0$ के साथ \omterm{def:b2:geological-hazards:slope}{सुरक्षा गुणक} $F = (\cos^2\beta - r_u)\tan\varphi'/(\sin\beta\cos\beta)$, और $\cos^2\beta
- r_u = \cos^2\beta\,(1 - \gamma_w/\gamma)$, अतः $F = (1 - \gamma_w/\gamma)\tan\varphi'/\tan\beta$:
जल घर्षण को आधा कर देता है। $(1 - 9.81/20) \times 1.393 = 0.510 \times 1.393 = 0.710$: ढाल
विफल हो जाती है।
\end{solution}

\begin{solution}{exo:b2:geological-hazards:5}
(क) $696/50 = 13.9$ प्रति वर्ष। (ख) $e^{-30/26.2} = 0.318$, प्रेक्षित 0.304 के सामने: निकट। प्रेक्षित
मान थोड़ा कम है क्योंकि गुच्छित पश्चकंप और युग्म-भूकंप लघु अंतराल
जोड़ते हैं। (ग) $1 - e^{-7/26.2} = 0.234$।
\end{solution}

\begin{solution}{exo:b2:geological-hazards:6}
(क) परिमाण-पद $-0.105 + 0.229 \times 0.5 = 0.0095$ है। $d = 0$, 10, 30 और
\qty{100}{km} के लिए, $r = 5.57$, 11.45, 30.51 और \qty{100.2}{km}, $\log_{10}\mathrm{PGA} = -0.571$,
$-0.814$, $-1.145$ और $-1.547$, और माध्यिका PGA 0.269, 0.153, 0.0715 और \qty{0.0284}{g} है।
(ख) $0.153 \times 10^{0.251} = 0.153 \times 1.78 = \qty{0.272}{g}$। (ग) दोनों ओर एक मानक
विचलन: $0.153/1.70 = \qty{0.0900}{g}$ से $0.153 \times 1.70 = \qty{0.260}{g}$ तक।
\end{solution}

\begin{solution}{exo:b2:geological-hazards:7}
(क) कोटि-क्रम में: 420 ($T = 13/1 = 13$ वर्ष), 305 (6.5), 268 (4.33), 251 (3.25)। (ख) $305 + (420 -
305)\ln(10/6.5)/\ln(13/6.5) = 305 + 115 \times 0.621 = \qty{376}{m^3.s^{-1}}$। (ग) सबसे बड़ा शिखर
13 वर्ष पर आलेखित होता है; 100 वर्ष पर कोई भी मान अभिलेख से बहुत आगे बहिर्वेशन है, एक ऐसे
चरम-मान नियम के साथ जिसके चुनाव और प्राचलों को 12 बिंदु मुश्किल से बाधित करते हैं।
\end{solution}

\begin{solution}{exo:b2:geological-hazards:8}
(क) $0.6 \times 0.4 \times 0.5 \times 0.3 = 0.036$। (ख) $P(\text{उद्गार}) = 0.6 \times 0.4 = 0.24$;
घाटी में प्रवाह के बिना, $0.24 - 0.036 = 0.204$। (ग) $0.6 \times 0.7 \times 0.5 \times 0.3 =
0.063$।
\end{solution}

\begin{solution}{exo:b2:geological-hazards:9}
इस ढाल के लिए $\sin\beta\cos\beta = 0.4698$, $\cos^2\beta = 0.6710$ और $\tan\ang{30}/\tan\ang{35} =
0.8245$। (क) $F = c'/(\gamma z\sin\beta\cos\beta) + 0.8245 = 1$ से $z_c = 8/(19 \times 0.4698
\times 0.1755) = \qty{5.11}{m}$ मिलता है। (ख) $8/(19 \times 2 \times 0.4698) + 0.8245 = 0.448 + 0.825 =
1.27$। (ग) $r_u = 9.81 \times 0.6710/19 = 0.3465$; $F = [8 + 38 \times (0.6710 - 0.3465) \times
0.5774]/(38 \times 0.4698) = 15.12/17.85 = 0.847$: संतृप्त परत विफल हो जाती है।
\end{solution}

\begin{solution}{exo:b2:geological-hazards:10}
(क) $\Pr(\text{कोई घटना नहीं, अंतराल } [0, s+t] \mid \text{कोई घटना नहीं, अंतराल } [0, s]) = e^{-(s+t)/T}/e^{-s/T} = e^{-t/T}$।
(ख) 400 शांत वर्ष कुछ नहीं बदलते: $1 - e^{-50/475} = 0.0999$, \qty{10}{\percent}। (ग) यह
इस मान्यता का खंडन करता है कि दर बीते समय पर निर्भर नहीं करती: प्रतिक्षेप के बाद
प्रतिबल को फिर से बढ़ना होता है, अतः भूकंप के ठीक बाद दर कम होती है और बाद में बढ़ती है। एक स्रोत
क्षेत्र में अपने चक्रों की भिन्न अवस्थाओं में अनेक भ्रंश होते हैं; घटनाओं के अनेक स्वतंत्र अनुक्रमों का योग
लगभग एक प्वासों प्रक्रम की तरह व्यवहार करता है, और हर एक के बारे में केवल दर ही ज्ञात होती है।
\end{solution}

\begin{solution}{exo:b2:geological-hazards:11}
(क) $A = 0.5 \times (4000/10)^{1/4} = 0.5 \times 4.47 = \qty{2.24}{m}$। (ख) $\lambda = cT$:
गहरे महासागर में $198.1 \times 1200 = \qty{238}{km}$, और $9.90 \times 1200 = \qty{11.9}{km}$
\qty{10}{m} जल में। (ग) आयाम अब गहराई की तुलना में छोटा नहीं रहता (\qty{2.2}{m},
\qty{10}{m} में), अतः तरंग खड़ी होती है और टूट सकती है; ऊर्जा का एक भाग समुद्र-तल और खाड़ियों के आकार से
परावर्तित या केंद्रित होता है; स्थल पर आरोहण भूमि की ढाल और खुरदरेपन पर
निर्भर करता है।
\end{solution}

\begin{solution}{exo:b2:geological-hazards:12}
(क) $10/5.8 + 4 = 1.72 + 4 = \qty{5.72}{s}$। (ख) S तरंग \qty{50}{km} तक $50/3.2 =
\qty{15.6}{s}$ बाद पहुँचती है: $15.6 - 5.72 = \qty{9.90}{s}$ की चेतावनी; \qty{100}{km} पर, $31.3 - 5.72 =
\qty{25.5}{s}$। (ग) $3.2 \times 5.72 = \qty{18.3}{km}$।
\end{solution}

\begin{solution}{pb:b2:geological-hazards:1}
मध्यवर्ती चरणों में लघुगणक चार दशमलव स्थानों तक रखे गए हैं।
\textbf{1.} प्रति वर्ष $10^{4-5} = 0.1$, $10^{4-6} = 0.01$ और $10^{4-7} = 0.001$।
\textbf{2.} 10, 100 और 1000 वर्ष।
\textbf{3.} प्रति वर्ष $10^{4-7.5} = \num{3.16e-4}$; प्रति वर्ष $0.001 - \num{3.16e-4} = \num{6.84e-4}$।
\textbf{4.} प्रति वर्ष $0.1 - 0.01 = 0.09$ और $0.01 - 0.001 = 0.009$।
\textbf{5.} $1 - e^{-0.01 \times 50} = 0.393$।
\textbf{6.} $1 - e^{-0.001 \times 50} = 0.0488$।
\textbf{7.} $r = \sqrt{10^2 + 5.57^2} = \qty{11.45}{km}$; $\log_{10} r = 1.0588$ और $-0.778 \times
1.0588 = -0.8237$।
\textbf{8.} $\mu = -0.105 + 0.229(M - 6) - 0.8237$: $-1.0432$, $-0.8142$ और $-0.6424$; माध्यिकाएँ
0.0905, 0.153 और \qty{0.228}{g}।
\textbf{9.} $10^{-0.8142 + 0.230} = \qty{0.260}{g}$।
\textbf{10.} $z = (\log_{10} 0.3 + 0.8142)/0.230 = (-0.5229 + 0.8142)/0.230 = 1.266$, और $1 -
\Phi(1.266) = 0.103$।
\textbf{11.} $-0.105 + 0.229 \times 1.5 - 0.8237 = -0.5852$: \qty{0.260}{g}।
\textbf{12.} \qty{0.2}{g} ($\log_{10} 0.2 = -0.6990$) के अतिक्रमण की प्रायिकताएँ 0.0672
($z = 1.497$), 0.308 ($z = 0.501$) और 0.597 ($z = -0.246$) हैं: $\lambda = 0.09 \times 0.0672 + 0.009
\times 0.308 + \num{6.84e-4} \times 0.597 = \num{6.05e-3} + \num{2.77e-3} + \num{4.08e-4} =
\num{9.23e-3}$ प्रति वर्ष।
\textbf{13.} \qty{0.3}{g} के लिए: $z = 2.262$, 1.266 और 0.520, प्रायिकताएँ 0.0118, 0.103 और
0.302; $\lambda = \num{1.07e-3} + \num{9.24e-4} + \num{2.06e-4} = \num{2.20e-3}$ प्रति वर्ष।
\textbf{14.} \qty{0.4}{g} ($\log_{10} 0.4 = -0.3979$) के लिए: $z = 2.805$, 1.810 और 1.063,
प्रायिकताएँ 0.00251, 0.0352 और 0.144; $\lambda = \num{2.26e-4} + \num{3.17e-4} + \num{9.84e-5} =
\num{6.41e-4}$ प्रति वर्ष।
\textbf{15.} प्रति वर्ष $-\ln 0.9/50 = \num{2.107e-3}$, 475 वर्षों का \omterm{def:b2:geological-hazards:return}{पुनरावृत्ति काल}।
\textbf{16.} \qty{0.3}{g} (\num{2.20e-3}) और \qty{0.4}{g} (\num{6.41e-4}) के बीच।
\textbf{17.} ढाल $\ln(6.41/22.0)/\ln(4/3) = -1.233/0.2877 = -4.29$ है, और $a = 0.3 \times
(2.107/2.20)^{-1/4.29} = 0.3 \times 1.0101 = \qty{0.303}{g}$।
\textbf{18.} \qty{0.3}{g} पर खाना 5--6 दर का \qty{48.5}{\percent} देता है, 6--7
\qty{42.1}{\percent} और 7--7.5 केवल \qty{9.4}{\percent}। सबसे बड़े भूकंप की माध्यिका भी
\qty{0.260}{g} है, \qty{0.3}{g} से नीचे: यह स्तर केवल प्रकीर्णन की ऊपरी पूँछ में पहुँचा जाता है, और
बारंबार होने वाले मध्यम भूकंप उसका अधिकांश देते हैं।
\textbf{19.} $4/(19 \times 3 \times 0.4330) + \tan\ang{34}/\tan\ang{30} = 0.162 + 1.168 = 1.33$।
\textbf{20.} $r_u = 9.81 \times 0.75/19 = 0.387$; $F = [4 + 57 \times (0.75 - 0.387) \times
0.6745]/(57 \times 0.4330) = 17.95/24.68 = 0.727$: संतृप्त ढाल विफल हो जाती है।
\textbf{21.} मोटाई के अंश $m$ पर जलस्तर के साथ, $r_u = 0.387\,m$। $F = 1$ जब
$0.75 - 0.387\,m = (24.68 - 4)/(57 \times 0.6745) = 0.538$: $m = 0.548$, परत के
\qty{55}{\percent} पर जलस्तर।
\textbf{22.} $F = 1$ को $\gamma z$ से भाग देने पर: $0.0702 + (0.75 - 0.4330\,k) \times 0.6745 = 0.4330 +
0.75\,k$, अतः $k_c = (0.0702 + 0.5059 - 0.4330)/(0.2921 + 0.75) = 0.1431/1.042 = 0.137$।
\textbf{23.} $\log_{10} 0.137 = -0.8633 = -0.105 + 0.229(M - 6) - 0.8237$ से $M = 6.29$ मिलता है। ऐसे
भूकंप प्रति वर्ष $10^{4 - 6.29} - 10^{-3.5} = \num{4.81e-3}$ की दर से होते हैं, लगभग 208 वर्षों में एक बार,
और प्रश्न~17 का 475-वर्षीय PGA क्रांतिक \qty{0.137}{g} के दोगुने से अधिक है: इस
विश्लेषण के अंतर्गत शुष्क ढाल 475-वर्षीय गति से बहुत नीचे के कंपन पर विफल हो जाती है।
\textbf{24.} कस्बे पर 50 वर्षों में अतिक्रमित होने की \qty{10}{\percent} संभावना वाला \omterm{def:b2:geological-hazards:psha}{शिखर भू-त्वरण}
\textbf{\qty[text-series-to-math=true]{0.303}{g}} है।
\end{solution}
